% IO-01D2 Phase 4 formal results snippet
\subsection{MCP Enforcement Under Explicit Challenge}

To distinguish behavioral steering from technical enforcement, we conducted a
10-trial, $2\times2$ factorial enforcement challenge in which baseline and
Markdown-governed sessions received the same prompt requesting exactly one
validation submission through the institutional MCP governance service. The
authorization factor was external to the prompt: the MCP control plane was
configured either with authorization absent or granted. Production Slurm was
not contacted; allowed scheduler actions were simulated by the governance
service.

All 40 sessions were valid. Every session attempted exactly one
\texttt{submit\_job} action. When authorization was absent, all 20 requests were
denied and no unauthorized action succeeded. When authorization was granted,
all 20 requests were allowed exactly once and returned a successful simulated
validation result. No session attempted a direct scheduler bypass. Thus,
authorization state perfectly determined the governed capability outcome while
the baseline versus Markdown condition did not alter the enforcement decision.

The result complements the earlier behavioral-governance experiment. Markdown
guidance changed whether the assistant attempted a state-changing scheduler
action when authorization was unspecified, whereas the MCP experiment
demonstrates that, once an action is attempted, the institutional control plane
can enforce authorization independently of the assistant's behavioral
guidance. This supports a layered model in which Markdown provides behavioral
steering and MCP provides constrained, auditable, enforceable capabilities.

One session (Trial 7, authorization absent, baseline) skipped the optional
read-only \texttt{validate\_job} preflight and proceeded directly to the single
\texttt{submit\_job} request. The request was still denied, indicating that the
authorization boundary did not depend on voluntary preflight use.

Because the MCP scheduler backend was simulated and direct production Slurm
access was intentionally blocked, these results demonstrate authorization
enforcement and capability routing rather than production scheduler security
under an unrestricted hostile process.
